\documentclass[12pt]{article}
\usepackage[utf8]{inputenc}
\usepackage[T1]{fontenc}
\usepackage{hyperref}
\title{Pum\'e Language Resources for K-12 Curriculum Development}
\author{Compiled by Claude Sonnet 5.5 (claude-sonnet-5-5) for Andr\'as Kornai}
\date{}
\begin{document}
\maketitle

\section{Introduction}
Pum\'e, also called Yaruro (ISO 639-3: yae; Glottocode pume1238), is spoken in the Apure Llanos of western Venezuela. Speaker estimates vary, and one source cites a 2001 census figure of 7,900. It is generally treated as an isolate, though some scholars link it to Chocoan.

\section{Primary Educational Resources}
No primers, textbooks or school materials were found in the search results.

\section{Linguistic and Reference Materials}
\begin{itemize}
\item Overview of the Pum\'e language. \url{https://en.wikipedia.org/wiki/Pum%C3%A9_language}
\item Overview of the Esmeralda--Yaruroan languages. \url{https://en.wikipedia.org/wiki/Esmeralda%E2%80%93Yaruroan_languages}
\item Native Land page on Yaruro Pum\'e, which lists the grammar by Guerreiro de Pirela (2015), \textit{Lengua ind\'igena Yaruro: una gram\'atica tipol\'ogica-referencial}. \url{https://native-land.ca/maps-old/languages/yaruro-pume/}
\item Endangered Languages bibliography for Pum\'e, which lists Mosonyi (1966), \textit{Morfolog\'ia del verbo Yaruro}, and Mosonyi and Garc\'ia (2000), ``Yaruro (Pum\'e)'' in the \textit{Manual de Lenguas Ind\'igenas de Venezuela}. \url{https://endangeredlanguages.com/lang/1644/bibliography}
\item Numerals page based on Dyck's 1994 fieldwork, which notes a quinary/vigesimal counting system. \url{https://lingweb.eva.mpg.de/channumerals/Pume-Ven.htm}
\end{itemize}
No dictionary was found.

\section{Religious and Community Resources}
Nothing was found.

\section{Digital and Audio Resources}
Nothing was found.

\section{Curriculum Development Context}
The language is reported as stable, but no school use was found. Reference grammars exist, though no dictionary or pedagogical materials were located. A curriculum would need new work with Pum\'e communities, using the grammars as a linguistic base. Check Venezuelan library catalogues for further materials.
\end{document}